\chapter{Introduction}\label{cap:introduction}
\glsresetall

% =====================================================================
% SKELETON. Target: 5--7 pages.
%
% Narrative arc for this chapter, in order:
%   1. Controllable facial expression synthesis matters (applications:
%      affective computing data, avatars, clinical/behavioural stimuli,
%      training data for FER).
%   2. Diffusion models made photorealistic face generation routine, but
%      control over *how much* of an expression is present is still coarse
%      (text prompts are categorical: "smiling" / "not smiling").
%   3. Prior art conditions on continuous AU intensity (\citet{varanka2024fineface})
%      but generates each intensity level as an INDEPENDENT still: nothing
%      binds identity or scene across the levels, and intensity control is
%      reported qualitatively.
%   4. This work reframes expression intensity as a SEQUENCE-level generation
%      parameter: one text prompt plus a per-frame intensity trajectory
%      yields a single temporally coherent clip.
%   5. What was measured so far, honestly (including the negative
%      calibration result) -> forward-reference to Cap05.
%
% Style: impersonal voice, no em-dashes, define every acronym on first use
% via \gls{}. Keep the detailed method out of here (it belongs in Cap04).
% =====================================================================

% =====================================================================
\section{Motivation}\label{sec:motivation}

Controllable synthesis of facial expressions has broad use, and in several of its applications
what matters is not whether a face is smiling but \emph{how much}. Affective-computing and
behavioural studies require stimuli whose expression intensity is varied in a controlled way, and
that intensity is not an incidental parameter of the stimulus. Spontaneous expressions in real
social interaction are frequently of low intensity, whereas conventional stimulus sets portray
emotions only at their peak, a mismatch that motivated the construction of graded-intensity sets
such as the \gls{adfes} Bath Intensity Variations \citep{wingenbach2016adfesbiv}. Its validation
also shows that the axis is behaviourally consequential: recognition accuracy rises monotonically
with intensity, from $56\%$ for low-intensity expressions to $68\%$ at intermediate and $75\%$ at
high intensity. The level at which an expression is rendered therefore determines what a study
using it can measure.

The same axis is a bottleneck for data. Deep \gls{fer} is limited above all by a shortage of
training data \citep{li2022fersurvey}, and intensity-annotated data is the scarcest kind, because
annotating it is slow: coding action-unit intensity under the \gls{facs} \citep{ekman1978facs} can
take an expert more than an hour for a single second of video \citep{tran2017deepcoder}. The
corpora that do carry intensity annotation quantise it coarsely, on a six-point discrete scale in
\gls{disfa} \citep{mavadati2013disfa} or three levels per clip in \gls{mead}
\citep{wang2020mead}. A generator able to render a specified, continuously graded intensity on
demand would relieve both the stimulus and the data problem, provided that the intensity it
renders can be verified properly.

Diffusion models have made photorealistic face generation routine, and text prompts give
convenient high-level control. That control, however, is categorical: a prompt distinguishes
``smiling'' from ``not smiling'', and although language can qualify degree, a phrase such as
``a slight smile'' carries no reproducible mapping onto a measurable intensity. The same words
render differently across prompts, seeds and models, and nothing in the interface allows a smile
to be commanded at a stated fraction of its full extent. What is needed is a numeric axis of
expression intensity on which a diffusion backbone can be conditioned directly, and, for the
sequences that behavioural stimuli and augmentation data require, an axis that varies across the
frames of a single coherent clip rather than across unrelated images.

% =====================================================================
\section{Problem Statement}\label{sec:problem}

% The per-paper analysis of the four closest systems lives in
% Chapter~\ref{cap:related-work} (source material:
% docs/experiments/2026-07-11-related-work-study.md). This section deliberately
% does NOT repeat it. It states the requirements a solution must satisfy, so
% that the contributions of Section~\ref{sec:contributions} and the evaluation
% protocol of Chapter~\ref{cap:method} can be read against them.

The problem is defined by what a solution would have to do. Three requirements follow from the
applications of Section~\ref{sec:motivation}, and a solution must satisfy them at the same time.

\begin{itemize}
    \item \textbf{Numeric Intensity Signal.} The intensity of the expression
    must be specified as a number instead of vague words, and that number must govern each frame
    individually, so that a whole trajectory from neutral to apex can be requested in a single
    instruction.

    \item \textbf{Identity and scene consistency.} If each intensity level is rendered as a separately
    generated image, any difference observed between two levels confounds the change in intensity with a change of identity, pose and scene. For the behavioral stimuli of
    Section~\ref{sec:motivation}, that confound is disqualifying: a response cannot be attributed
    to the intensity of an expression if the face displaying it also changed.

    \item \textbf{Independent Verification.} What was commanded must be checkable against
    what was rendered, and the instrument performing the check must not be the one that produced
    the training labels. Without an independent check the commanded intensity cannot be properly verified, and a model can score well by reproducing the biases of its
    own annotator.
\end{itemize}

No existing system satisfies the three simultaneously. Chapter~\ref{cap:related-work} examines
the closest ones and shows which requirement each relinquishes: expression transfer accepts no
numeric intensity at all; editing methods take a source image as their reference, and so command
a change rather than a value; and action-unit-conditioned generation renders each level as an
independent image.

% =====================================================================
\section{Research Questions}\label{sec:research-questions}

% Keep to two or three. Each must be answerable by an experiment in Cap05
% or by the planned work in Cap06.

\noindent\textbf{Task formulation.} Let a \emph{clip} denote a sequence of $F$ frames
that share one identity and scene. Given a text prompt and a per-frame intensity
trajectory $(s_1,\dots,s_F)$ with each $s_t \in [0,1]$, the objective is to generate a
single clip in which the detected intensity of a given action unit in frame $t$ tracks
$s_t$. This work instantiates the axis with AU12 (lip-corner puller), a proxy for smile
intensity. 

Prior conditional-generation methods instead render each intensity level as an independent image. The three research questions below ask, in turn, whether this trajectory-following is achievable on a frozen \gls{t2v} backbone (RQ1), whether generating the trajectory jointly as one sequence benefits cross-level identity (RQ2), and whether the resulting control is absolute or relative to the commanded trajectory (RQ3). This formulation is restated and made precise
in Chapter~\ref{cap:method}, Section~\ref{sec:method-overview}.

\medskip
\noindent\textbf{RQ1}: Can a scalar expression-intensity axis be injected into a frozen
\gls{t2v} backbone by the same conditioning mechanism that \citet{yuan2024genphoto} apply
to camera parameters, such that the detected AU12 intensity follows a commanded per-frame
trajectory?

\noindent\textbf{RQ2}: Does generating the intensity levels jointly, as a single
sequence, preserve identity across levels better than generating each level as an
independent still?

\noindent\textbf{RQ3}: Is the resulting control absolute, so that a commanded value maps
to a fixed detected intensity, or is it relative to the commanded trajectory, and what
property of the training data determines which?

% =====================================================================
\section{Objectives}\label{sec:objectives}

\subsection{General Objective}\label{subsec:general-objective}

The general objective of this work is to adapt the camera-parameter conditioning mechanism of a
frozen text-to-video diffusion backbone into a facial-expression intensity axis, so that a text
prompt and a per-frame intensity trajectory generate a single identity-consistent clip whose
detected expression intensity follows the command, and to evaluate that control with a protocol
that separately measures trajectory following, absolute calibration, and cross-frame identity.

\subsection{Specific Objectives}\label{subsec:specific-objectives}

\begin{itemize}
    \item Adapt the camera-parameter conditioning mechanism of \citet{yuan2024genphoto} to a facial-expression intensity axis.
    \item Build continuous per-frame AU12 supervision from video (\gls{mead}).
    \item Define an evaluation protocol measuring ramp following, absolute calibration, and cross-frame identity.
    \item Compare head-to-head against the closest prior art under matched prompts.
    \item Characterise the failure modes: calibration, demographic bias, and non-monotonic trajectories.
\end{itemize}

% =====================================================================
\section{Hypotheses}\label{sec:hypotheses}

% Verdicts deliberately kept OUT of this chapter: the hypotheses are stated
% here and adjudicated in Cap05, one subsection each. H3 is refuted there by
% the 2026-07-14 calibration experiment; keeping a refuted hypothesis and
% reporting the refutation is stronger than quietly deleting it.
%   H1 -> Section~\ref{sec:res-ramp}
%   H2 -> Section~\ref{sec:res-headtohead}
%   H3 -> Section~\ref{sec:res-calibration}

Three hypotheses follow from the research questions above. Each is adjudicated against the
measurements of Chapter~\ref{cap:preliminary-results}.

\begin{itemize}
    \item \textbf{H1}: Per-frame scalar conditioning on a frozen backbone yields a detected
    intensity that follows the commanded trajectory.
    \item \textbf{H2}: Joint sequence generation preserves identity across intensity levels
    better than independent per-level generation.
    \item \textbf{H3}: The control is absolute, that is, a commanded value maps to the same
    detected intensity regardless of the surrounding trajectory. Conditioning applied per frame
    and supervised by a per-frame label should in principle fix a mapping from commanded value to
    rendered intensity that is independent of the neighbouring frames.
\end{itemize}

% =====================================================================
\section{Contributions}\label{sec:contributions}

% Claim only what survives the FineFace head-to-head
% (docs/experiments/2026-07-21-fineface-headtohead.md):
%   - ramp-following is a TIE with FineFace, not a win;
%   - FineFace BEATS this work on absolute calibration;
%   - this work wins cross-frame identity, modestly.
% The defensible contribution is therefore the sequence-native formulation
% plus the measurement protocol, NOT better expression control. Any
% "first AU-conditioned generation" phrasing is dead: FineFace (July 2024)
% has primacy. Also avoid "fine-grained" as a headline term, since three of
% the four related papers use it with different meanings.

\noindent This work makes three contributions.

\begin{itemize}
    \item \textbf{C1 (Formulation).} Expression intensity is cast as a \emph{sequence-level}
    generation parameter: one text prompt and one per-frame trajectory yield a single
    identity-consistent clip, whereas the closest prior art \citep{varanka2024fineface} (FineFace)
    produces $N$ independent stills bound only by a shared random seed. Delivered by the
    method overview and conditioning design and evaluated in the head-to-head comparison in Section~\ref{sec:res-headtohead}.
    \item \textbf{C2 (Evaluation protocol).} A protocol that separates three axes commonly
    conflated, namely trajectory-following, absolute calibration, and cross-frame
    identity, and that guards the label/evaluation circularity shared by the closest works
    by scoring with an estimator that produced no training labels. Delivered by
    Section~\ref{sec:method-evaluation} and applied across
    Sections~\ref{sec:res-ramp}--\ref{sec:res-headtohead}.
    \item \textbf{C3 (Empirical characterization).} A quantified commanded-versus-detected
    calibration curve, reported as a single comparable scalar for both this work and FineFace,
    which reports intensity control only qualitatively. The measurements show that ramp-only video
    supervision produces control that is relative to the commanded trajectory rather than absolute. 
    Delivered by Section~\ref{sec:res-calibration}.
\end{itemize}

\noindent Together these establish the contribution as the sequence-native formulation and the evaluation protocol, while also having comparable results with state-of-the-art models.

% =====================================================================
\section{Document Organization}\label{sec:organization}

The remainder of this document is organised as follows. Chapter~\ref{cap:background} introduces
the diffusion, conditioning, temporal-generation, and facial-action-coding machinery on which the
method rests. Chapter~\ref{cap:related-work} positions the work against the closest expression
transfer, editing, and generation systems. Chapter~\ref{cap:method} details the proposed method:
the frozen backbone, the expression conditioning, the training data and setup, and the evaluation
protocol. Chapter~\ref{cap:preliminary-results} reports what has been measured so far. Finally, Chapter~\ref{cap:schedule} lists the completed and remaining
work with the schedule to completion.
